\chapter*{はじめに}
\addcontentsline{toc}{chapter}{はじめに}

本ノートは，小平邦彦『複素多様体と複素構造の変形 I』(1968) の輪読ゼミにおいて
作成した記録である．本では省略されている導出や細部を自分の手で補うことを主な
目的とする．

\section*{ゼミノートの方針}

\begin{enumerate}
  \item テキストに書かれていることを証明する．主張を自分の言葉で再構成する．
  \item 記号など文脈上明らかなものは適宜省略する．
  \item 理解を優先する：省略されている，あるいはごまかされている箇所はすべて
        説明する．
  \item 証明は自分で再構成する，あるいは議論の流れを必要に応じて修正する．
  \item 疑問点や補足はできるだけ多く書き留める（予習ノートとは別に記録する）．
  \item 発表ノートを作る（理想的には，ノートなしでも発表できるように準備する）．
\end{enumerate}
